\chapter{Privilegien}
\label{ch:players/privileges}

Privilegien,\index{privileg@Privileg} kurz \emph{Privs}, erlauben
Spielern bestimmte Aktionen. Serverbetreiber können Privilegien vergeben
und entziehen und so steuern, was jeder Spieler darf.

\section{Wann sind Privilegien
sinnvoll?}\label{wann-sind-privilegien-sinnvoll}

Ein Privileg sollte einem Spieler erlauben, etwas zu \emph{tun}.
Privilegien sind \textbf{nicht} dazu da, eine Klasse oder einen Rang
anzuzeigen.

\textbf{Gute Privilegien:}

\begin{itemize}
\tightlist
\item
  \texttt{interact}
\item
  \texttt{shout}
\item
  \texttt{noclip}
\item
  \texttt{fly}
\item
  \texttt{kick}
\item
  \texttt{ban}
\item
  \texttt{vote}
\item
  \texttt{worldedit}
\item
  \texttt{area\_\allowbreak{}admin} -- Admin-Funktionen eines einzelnen
  Mods sind in Ordnung
\end{itemize}

\textbf{Schlechte Privilegien:}

\begin{itemize}
\tightlist
\item
  \texttt{moderator}
\item
  \texttt{admin}
\item
  \texttt{elf}
\item
  \texttt{dwarf}
\end{itemize}

\section{Privilegien deklarieren}\label{privilegien-deklarieren}

Ein neues Privileg legst du mit
\texttt{register\_\allowbreak{}privilege} an:

\begin{Shaded}
\begin{Highlighting}[]
\VariableTok{core}\OperatorTok{.}\NormalTok{register\_privilege}\OperatorTok{(}\StringTok{"vote"}\OperatorTok{,} \OperatorTok{\{}
    \VariableTok{description} \OperatorTok{=} \StringTok{"Darf abstimmen"}\OperatorTok{,}
    \VariableTok{give\_to\_singleplayer} \OperatorTok{=} \KeywordTok{true}\OperatorTok{,}
\OperatorTok{\})}
\end{Highlighting}
\end{Shaded}

\texttt{give\_\allowbreak{}to\_\allowbreak{}singleplayer} ist
standardmäßig \texttt{true}; in der Definition oben wäre es also
eigentlich nicht nötig.

\section{Privilegien prüfen}\label{privilegien-pruxfcfen}

So prüfst du schnell, ob ein Spieler alle nötigen Privilegien hat:

\begin{Shaded}
\begin{Highlighting}[]
\KeywordTok{local} \VariableTok{has}\OperatorTok{,} \VariableTok{missing} \OperatorTok{=}
    \VariableTok{core}\OperatorTok{.}\NormalTok{check\_player\_privs}\OperatorTok{(}\VariableTok{player\_or\_name}\OperatorTok{,} \OperatorTok{\{}
        \VariableTok{interact} \OperatorTok{=} \KeywordTok{true}\OperatorTok{,}
        \VariableTok{vote} \OperatorTok{=} \KeywordTok{true}\OperatorTok{,}
    \OperatorTok{\})}
\end{Highlighting}
\end{Shaded}

Hier ist \texttt{has} gleich \texttt{true}, wenn der Spieler alle
nötigen Privilegien hat. Ist \texttt{has} gleich \texttt{false}, enthält
\texttt{missing} eine Schlüssel-Wert-Tabelle der fehlenden Privilegien.

\begin{Shaded}
\begin{Highlighting}[]
\KeywordTok{local} \VariableTok{has}\OperatorTok{,} \VariableTok{missing} \OperatorTok{=} \VariableTok{core}\OperatorTok{.}\NormalTok{check\_player\_privs}\OperatorTok{(}\VariableTok{name}\OperatorTok{,} \OperatorTok{\{}
    \VariableTok{interact} \OperatorTok{=} \KeywordTok{true}\OperatorTok{,}
    \VariableTok{vote} \OperatorTok{=} \KeywordTok{true}\OperatorTok{,}
\OperatorTok{\})}

\ControlFlowTok{if} \VariableTok{has} \ControlFlowTok{then}
    \FunctionTok{print}\OperatorTok{(}\StringTok{"Spieler hat alle Privilegien!"}\OperatorTok{)}
\ControlFlowTok{else}
    \FunctionTok{print}\OperatorTok{(}\StringTok{"Spieler fehlen Privilegien: "} \OperatorTok{..} \FunctionTok{dump}\OperatorTok{(}\VariableTok{missing}\OperatorTok{))}
\ControlFlowTok{end}
\end{Highlighting}
\end{Shaded}

Brauchst du die fehlenden Privilegien nicht, kannst du
\texttt{check\_\allowbreak{}player\_\allowbreak{}privs} direkt in die
\texttt{if}-Anweisung schreiben.

\begin{Shaded}
\begin{Highlighting}[]
\ControlFlowTok{if}
    \KeywordTok{not} \VariableTok{core}\OperatorTok{.}\NormalTok{check\_player\_privs}\OperatorTok{(}\VariableTok{name}\OperatorTok{,} \OperatorTok{\{} \VariableTok{interact} \OperatorTok{=} \KeywordTok{true} \OperatorTok{\})}
\ControlFlowTok{then}
    \ControlFlowTok{return} \KeywordTok{false}\OperatorTok{,} \StringTok{"Dafür brauchst du interact!"}
\ControlFlowTok{end}
\end{Highlighting}
\end{Shaded}

\section{Privilegien lesen und
setzen}\label{privilegien-lesen-und-setzen}

Die Privilegien eines Spielers kannst du lesen und ändern, egal ob er
online ist oder nicht.

\begin{Shaded}
\begin{Highlighting}[]
\KeywordTok{local} \VariableTok{privs} \OperatorTok{=} \VariableTok{core}\OperatorTok{.}\NormalTok{get\_player\_privs}\OperatorTok{(}\VariableTok{name}\OperatorTok{)}
\FunctionTok{print}\OperatorTok{(}\FunctionTok{dump}\OperatorTok{(}\VariableTok{privs}\OperatorTok{))}

\VariableTok{privs}\OperatorTok{.}\VariableTok{vote} \OperatorTok{=} \KeywordTok{true}
\VariableTok{core}\OperatorTok{.}\NormalTok{set\_player\_privs}\OperatorTok{(}\VariableTok{name}\OperatorTok{,} \VariableTok{privs}\OperatorTok{)}
\end{Highlighting}
\end{Shaded}

Privilegien werden immer als Schlüssel-Wert-Tabelle angegeben: Der
Schlüssel ist der Name des Privilegs, der Wert ein Wahrheitswert.

\begin{Shaded}
\begin{Highlighting}[]
\OperatorTok{\{}
    \VariableTok{fly} \OperatorTok{=} \KeywordTok{true}\OperatorTok{,}
    \VariableTok{interact} \OperatorTok{=} \KeywordTok{true}\OperatorTok{,}
    \VariableTok{shout} \OperatorTok{=} \KeywordTok{true}\OperatorTok{,}
\OperatorTok{\},}
\end{Highlighting}
\end{Shaded}

\section{Privilegien zu basic\_privs
hinzufügen}\label{privilegien-zu-basic_privs-hinzufuxfcgen}

Spieler mit dem Privileg \texttt{basic\_\allowbreak{}privs} dürfen eine
begrenzte Auswahl an Privilegien vergeben und entziehen. Üblicherweise
bekommen Moderatoren dieses Privileg, damit sie \texttt{interact} und
\texttt{shout} vergeben und entziehen können, sich selbst oder anderen
aber keine leichter missbrauchbaren Privilegien wie \texttt{give} oder
\texttt{server} geben können.

Welche Privilegien deine Moderatoren vergeben und entziehen dürfen,
legst du über die Einstellung \texttt{basic\_\allowbreak{}privs} fest.

Standardmäßig hat \texttt{basic\_\allowbreak{}privs} diesen Wert:

\begin{Shaded}
\begin{Highlighting}[]
\NormalTok{basic\_privs = interact, shout}
\end{Highlighting}
\end{Shaded}

Um \texttt{vote} hinzuzufügen, änderst du ihn so:

\begin{Shaded}
\begin{Highlighting}[]
\NormalTok{basic\_privs = interact, shout, vote}
\end{Highlighting}
\end{Shaded}

Damit dürfen Spieler mit \texttt{basic\_\allowbreak{}privs} das Privileg
\texttt{vote} vergeben und entziehen.
